\documentclass{article}
\usepackage{graphicx}
\usepackage{amsmath}
\usepackage{amsfonts}
\usepackage{bookmark}
\usepackage{booktabs}
\usepackage{exercises}
\graphicspath{ {} }
\usepackage{indentfirst}
\usepackage{amsthm}
\usepackage{yhmath}
\usepackage{geometry}
\usepackage{placeins}
\usepackage{csquotes}
\usepackage{pgfplots}
\newcommand*{\oldneg}{\mathord{\sim}}
\newcommand\showdiv[1]{\overline{\smash{)}#1}}
\newcommand\Mydiv[1]{%
\smash{\raise.3ex\hbox{$\big)$}}\mkern-8mu
        \overline{\enspace\strut#1}}
\setlength{\tabcolsep}{5pt}
\theoremstyle{definition}
\newtheorem{definition}{Definition}

\setlength{\parindent}{.5in}
\setlength{\parskip}{12pt}
\renewcommand{\arraystretch}{2}
\geometry{margin=1in}
\newcommand{\?}{\stackrel{?}{=}}
\pgfplotsset{width=7cm,compat=1.18}\usepgfplotslibrary{statistics}

\title{7.5 Answers}
\begin{document}
\maketitle

\section*{1}
The repeated composition of a function with itself is called iteration.

\section*{2}
An explicit rule for a sequence gives a rule for finding the $n$th tern of the sequence using only $n$. A recursive rule for a sequence gives a rule 
for finding the $n$th member of the sequence that could use $n$, but also any other previous terms.

\section*{10}
$2,4,2,-2,-4$ 

\section*{12}
A

\section*{14}
We know that this is a geometric sequence with $a_1=3$ and $d=4$. Thus we can model it with the following recursive rule:
$a_1=3$ and $a_n=a_{n-1}\times 4$.

\section*{19}
We know that this is an arithmetic sequence with $a_1=54$ and $d=-11$. Thus we can model it with the following recursive rule:
$a_1=54$ and $a_n=a_{n-1}-11$.

\section*{20}
This sequence can be modeled by: $a_1=3$, $a_2=5$ and $a_n=a_{n-1}\times a_{n-2}$.

\section*{21}
This sequence can be modeled by: $a_1=16$, $a_2=9$ and $a_n=a_{n-2}-a_{n-1}$.

\section*{31}
The first three iterates are $f_1(x_0)=f(x)$, $f(x_1)=f(f(x))$, and $f(x_2)=f(f(f(x)))$. $f(x_0)=f(1)=1^2-1+2=2$ meaning
$f(x_1)=f(2)=2^2-2+2=4$ and $f(x_2)=f(4)=4^2-4+2=14$.

\section*{32}
The first three iterates are $f_1(x_0)=f(x)$, $f(x_1)=f(f(x))$, and $f(x_2)=f(f(f(x)))$. $f(x_0)=f(2)=-3(2)^2+2(2)=-8$ meaning
$f(x_1)=f(-8)=-3(-8)^2+2(-8)=-208$ and $f(x_3)=-3(-208)^2+2(-208)=-130208$.

\section*{33}
C

\section*{36}
A recursive rule to model this sequence is $a_1=5$ and $a_n=a_{n-1}\times\sqrt{3}$.

\section*{38}
A recursive rule to model this sequence is $a_1=8$, $a_2=4$ and $a_n=\frac{a_{n-2}}{a_{n-1}}$.

\section*{39}
A recursive rule to model this sequence is $a_1=-3$, $a_2=-2$ and $a_n=-a_{n-2}-a_{n-1}$.

\section*{41}
Suppose we had a function $f$ with an initial value of $x_0$, and such that the functions first three iterates are $x_1=2,x_2=2$ and $x_3=8$. 
Now, $x_1=f(x_0)$, $x_2=f(x_1)$, and $x_3=f(x_2)$. Thus we would have $2=x_2=f(x_1)=f(2)$ and $8=x_3=f(x_2)=f(2)$. Thus we would have $f(2)=8$ and $f(2)=2$, a
contradiction. This means that this situation cannot occur.

\section*{42}
(a): We are given a recursive sequence defined by $a_1=3,a_2=10$ and $a_n=4+2a_{n-1}-a_{n-2}$. The first five terms of the sequence are
$3,10,4+2(10)-3=21,4+2(21)-10=36,4+2(36)-21=55$.


(b): Note that the differences are in the following sequence: $10-3=7,21-10=11,36-21=15,$ and $55-36=19$. Each term increases by a common difference of 4, meaning this is arithmetic. 
The sequence of the differences can therefore be expressed by $d_n=7+4(n-1)$.


(c): We have a sequence representing the differences in each successive term. What this means is that the $n$th term in the recursive series is equal to $3$ plus 
the sum of the differences up to that 
point, that is (because this is an arithmetic series): $s_n=n\frac{d_1+d_n}{2}=n\frac{7+7+4(n-1)}{2}=n\frac{14+4(n-1)}{2}=n(7+2(n-1))$. 
Thus an explicit rule for the sequence is $a_n=3+s_{n-1}=3+(n-1)(7+2(n-2))$.

\section*{47}
(This answer did not seem correct as the maintenance level was equal to the dosage level (which seems unexpected), but I could not see what went wrong)\newline
(a): A rule to model the amount of the drug in the bloodstream after $n$ doses is: $a_n=a_{n-1}\times 0.7+20$.


(b): In fact, we can write an explicit rule for the above sequence as well, which will help us determine the maintenence level. This can be done in the same manner as in Problem 63 of 
7.3:
\[ a_n=20\left(\dfrac{1-(\frac{1}{0.7})^{n-1}}{1-\frac{1}{0.7}}\right)(0.7)^{n-1}+20=20\left(\dfrac{1-(0.7)^{n-1}}{0.3}\right)(0.7)^{n-1}+20 \]
\[ =\frac{200}{3}(1-(0.7)^{n-1})(0.7)^{n-1}+20 \]
To determine what this value approaches over an extended period of time, we need to find its limit:
\[ \lim_{n\rightarrow\infty}\frac{200}{3}(1-(0.7)^{n-1})(0.7)^{n-1}+20 \]
\[ =\lim_{n\rightarrow\infty}\left(\frac{200}{3}(1-(0.7)^{n-1})(0.7)^{n-1}\right)+\lim_{n\rightarrow\infty} 20 \]
\[ =\lim_{n\rightarrow\infty}\left(\frac{200}{3}(1-(0.7)^{n-1})\right)\times\lim_{n\rightarrow\infty}(0.7)^{n-1}+20 \]
\[ =\lim_{n\rightarrow\infty}\left(\frac{200}{3}(1-(0.7)^{n-1})\right)\times 0+20 \]
\[ =20 \]


(c): If we double the dosage $d$, the maintenance level would be:
\[ \lim_{n\rightarrow\infty}\frac{2d}{0.3}(1-(0.7)^{n-1})(0.7)^{n-1}+2d \]
\[ =2d \]

\end{document}
